\section{Introduction}
\label{sec:intro}

Between July 2022 and September 2023 the ECB raised its policy rates by 450
basis points, the fastest tightening in the history of the euro. Real house
prices fell almost everywhere in the currency union, but by very different
amounts. One reason to expect such differences comes from mortgage
contracts. A Spanish or Finnish borrower, whose rate tracks the Euribor,
felt the tightening in the monthly payment within months. A German or
French borrower, whose rate was fixed years earlier, often felt nothing at
all. If monetary policy reaches household budgets at such different speeds,
the same decision in Frankfurt could move national housing markets by very
different amounts.

Whether it actually does is an open question. The euro area combines a
single monetary policy with mortgage systems that differ more than almost
any other institutional feature of its member states, and yet direct
evidence on house prices is scarce. A well established literature shows
that mortgage structure shapes the transmission of policy into household
consumption \citep{calza2013, dimaggio2017, cloyne2020}, and recent work
documents heterogeneous housing responses across countries and regions
\citep{corsetti2022, battistini2025}. Whether the fixed versus variable
character of a national mortgage market translates a common shock into
systematically different house price responses has, to my knowledge, not
been tested directly. As long as this evidence is missing, the common
presumption that variable rate countries carry more of the burden of a
tightening rests on indirect evidence, and neither the ECB nor national
macroprudential authorities can judge how unevenly a common decision lands
on national housing markets.

This paper asks the question directly. I build a quarterly panel of eleven
euro area countries from 1999 to 2025 and estimate panel local projections
\citep{jorda2005} of real house prices on high frequency monetary policy
surprises from the EA-MPD of \citet{altavilla2019}. Because raw surprises
mix policy news with central bank information about the economy, I clean
them with the sign restriction of \citet{jarocinski2020}. The cleaned shock
is then interacted with a binary classification of countries into fixed
rate dominant and variable rate dominant mortgage markets, following
\citet{badarinza2018}. Country fixed effects, lagged macroeconomic controls
and Driscoll Kraay standard errors complete the specification.

The short answer is no, although with important qualifications. Averaged
across countries, a
surprise tightening lowers real house prices somewhat over the first year,
though over the full sample this response is close to zero and imprecise.
Since 2009 the immediate response is negative and statistically
significant. On the central question, however, I find no robust evidence
that house prices in variable rate countries respond more strongly. The
estimated difference between the two country groups is small, it changes
sign depending on whether standard dynamic controls are included, and it
is not robustly significant anywhere. One further result stands out along
the way. Without the information correction, house prices appear to rise
after a tightening, so separating policy from information shocks is
necessary to get the housing results right.

Three possible objections should be addressed right away. The shock could be
mismeasured, but the series behaves as it should in well known ECB
episodes, and the contrast between raw and cleaned surprises matches what
the information effect literature predicts. The binary mortgage
classification is coarse, which is why I interpret the interaction as a
descriptive difference between country groups rather than a pure causal
effect of mortgage structure. Finally, eleven countries and one common
shock cannot deliver precise interaction estimates. That is true, and I
treat it openly as part of the answer. The confidence intervals show how
much heterogeneity the data can and cannot rule out.

The paper makes three contributions. It provides direct evidence on house
prices, rather than consumption or mortgage payments, for the mortgage
structure channel in the euro area, on a sample long enough to include the
2022 to 2023 tightening and the subsequent easing. It shows that the
information content of ECB announcements matters a great deal for
estimated housing responses. And it offers a transparent assessment of what a cross
country design can identify, which qualifies existing cross country
evidence and strengthens the case for regional data along the lines of
\citet{battistini2025}.
